% Soubory musí být v kódování, které je nastaveno v příkazu \usepackage[...]{inputenc}

\documentclass[%        Základní nastavení
    %  draft,    				  % Testovací překlad
    12pt,       				% Velikost základního písma je 12 bodů
    a4paper,    				% Formát papíru je A4
    oneside,      			% Jednostranný tisk
    %twoside,      			% Dvoustranný tisk - kapitoly začínají na lichých stranách
    electronic,   			% Elektronické PDF - symetrické okraje, hybridní prázdné stránky (pro tisk zakomentujte tuto volbu)
    unicode,						% Záložky a metainformace ve výsledném PDF budou v kódování unicode
]{report}				    	% Dokument třídy 'zpráva', vhodná pro sazbu závěrečných prací s kapitolami

\usepackage[utf8]		  %	Kódování zdrojových souborů je UTF-8
	{inputenc}					% Balíček pro nastavení kódování zdrojových souborů

\usepackage{sectsty}
	%přetypuje nadpisy všech úrovní na bezpatkové, kromě \chapter, která je přenastavena zvlášť v thesis.sty
	\allsectionsfont{\sffamily}

\usepackage{graphicx} % Balíček 'graphicx' pro vkládání obrázků
											% Nutné pro vložení logotypů školy a fakulty

\usepackage[          % Balíček 'acronym' pro sazby zkratek a symbolů
	nohyperlinks				% Nebudou tvořeny hypertextové odkazy do seznamu zkratek
]{acronym}						
											% Nutné pro použití prostředí 'acronym' balíčku 'thesis'

\usepackage[
	breaklinks=true,		% Hypertextové odkazy mohou obsahovat zalomení řádku
	hypertexnames=false % Názvy hypertext. odkazů budou tvořeny nezávisle na názvech TeXu
]{hyperref}						% Balíček 'hyperref' pro sazbu hypertextových odkazů
											% Nutné pro použití příkazu 'pdfsettings' balíčku 'thesis'

\usepackage{pdfpages} % Balíček umožňující vkládat stránky z PDF souborů
                      % Nutné při vkládání titulních listů a zadání přímo
                      % ve formátu PDF z informačního systému

\usepackage{enumitem} % Balíček pro nastavení mezerování v odrážkách
  \setlist{topsep=0pt,partopsep=0pt,noitemsep} % konkrétní nastavení

\usepackage{cmap} 		% Balíček cmap zajišťuje, že PDF vytvořené `pdflatexem' je
											% plně "prohledávatelné" a "kopírovatelné"

%\usepackage{upgreek}	% Balíček pro sazbu stojatých řeckých písmem
											%% např. stojaté pí: \uppi
											%% např. stojaté mí: \upmu (použitelné třeba v mikrometrech)
											%% pozor, grafická nekompatibilita s fonty typu Computer Modern!
                      
%\usepackage{amsmath} %balíček pro sabu náročnější matematiky                 

\usepackage{dirtree}	% sazba adresářové struktury
                      % vhodné pro prezentaci obsahu elektronické přílohy (např. CD)

\usepackage{algorithm2e} %balíček pro sazbu algoritmů/pseudokódu


%%%%%%%%%%%%%%% FONT VAFLE PRO TITULNÍ STRANU %%%%%%%%%%%%%%%
%% Font map file configuration for Vafle font
\pdfmapfile{=xlatex/font/vafle.map}

%% Definice pro kompatibilitu s LuaTeX
\ifcsname directlua\endcsname
  \protected\def\pdfmapfile{\pdfextension mapfile }
\fi
%%%%%%%%%%%%%%% FONT VAFLE PRO TITULNÍ STRANU %%%%%%%%%%%%%%%

\usepackage[formats]{listings}	% Balíček pro sazbu zdrojových textů
\lstset{              % nastavení
%	Definice jazyka použitého ve výpisech
%    language=[LaTeX]{TeX},	% LaTeX
%	language={Matlab},		% Matlab
	language={C},           % jazyk C
    basicstyle=\ttfamily,	% definice základního stylu písma
    tabsize=2,			% definice velikosti tabulátoru
    inputencoding=utf8,         % pro soubory uložené v kódování UTF-8
		columns=fixed,  %fixed nebo flexible,
		fontadjust=true %licovani sloupcu
    extendedchars=true,
    literate=%  definice symbolů s diakritikou
    {á}{{\'a}}1
    {č}{{\v{c}}}1
    {ď}{{\v{d}}}1
    {é}{{\'e}}1
    {ě}{{\v{e}}}1
    {í}{{\'i}}1
    {ň}{{\v{n}}}1
    {ó}{{\'o}}1
    {ř}{{\v{r}}}1
    {š}{{\v{s}}}1
    {ť}{{\v{t}}}1
    {ú}{{\'u}}1
    {ů}{{\r{u}}}1
    {ý}{{\'y}}1
    {ž}{{\v{z}}}1
    {Á}{{\'A}}1
    {Č}{{\v{C}}}1
    {Ď}{{\v{D}}}1
    {É}{{\'E}}1
    {Ě}{{\v{E}}}1
    {Í}{{\'I}}1
    {Ň}{{\v{N}}}1
    {Ó}{{\'O}}1
    {Ř}{{\v{R}}}1
    {Š}{{\v{S}}}1
    {Ť}{{\v{T}}}1
    {Ú}{{\'U}}1
    {Ů}{{\r{U}}}1
    {Ý}{{\'Y}}1
    {Ž}{{\v{Z}}}1
}

%%%%%%%%%%%%%%%%%%%%%%%%%%%%%%%%%%%%%%%%%%%%%%%%%%%%%%%%%%%%%%%%%
%%%%%%%%%%  Nastavení citačního stylu pomocí BIBLATEX  %%%%%%%%%%
%%%%%%%%%%%%%%%%%%%%%%%%%%%%%%%%%%%%%%%%%%%%%%%%%%%%%%%%%%%%%%%%%
%% Sekci zakomentujte pokud využíváte jen bibitem

\usepackage{csquotes}   % Balíček pro správné sazby uvozovek v citacích (doporučeno pro použití s biblatex)

\usepackage[
  backend=biber,        % typ zpracování
  style=iso-numeric,    % style složení citace
  defernumbers=true,    % reset číslování pro více literatur
  maxnames=5,           % Zobrazí max. počet autorů
  minnames=5,           % Když více, zobrazí et al.
  % urldate=iso,        % formát data přístupu podle ISO (YYYY-MM-DD)
]{biblatex}

% "a" před posledním autorem
\DefineBibliographyExtras{czech}{%
  \renewcommand*{\finalnamedelim}{\space a\space}
}

%% Formátování urldate podle českých norem (ISO 690). Aktuálně nastavuje špatně pozici před URL.
%% Doporučení: vypište text "[citováno YYYY-MM-DD]." do položky note
%% Zobrazí: [citováno YYYY-MM-DD].
% \DeclareFieldFormat{urldate}{\mkbibbrackets{citovánao\space#1}}

\addbibresource{text/literatura.bib} % Soubor s literaturou

\DeclareFieldFormat{labelnumberwidth}{[#1]} % Nastaví formát čísla v bibliography se zavorkama

\setlength{\biblabelsep}{1em}   % Nastaví mezeru za závorkou
\AtBeginBibliography{\small}    % velikost písma liteartury


%%%%%%%%%%%%%%%%%%%%%%%%%%%%%%%%%%%%%%%%%%%%%%%%%%%%%%%%%%%%%%%%%
%%%%%%%%%%%%%%%%%%%%%%%%%%%%%%%%%%%%%%%%%%%%%%%%%%%%%%%%%%%%%%%%%
%%%%%%%%%%%%%%%%%%  SEKCE PRO VLASTNÍ BALÍČKY  %%%%%%%%%%%%%%%%%%
%%%%%%%%%%%%%%%%%%%%%%%%%%%%%%%%%%%%%%%%%%%%%%%%%%%%%%%%%%%%%%%%%


%\usepackage[table,xcdraw]{xcolor}
%\usepackage{pifont}
%\usepackage{tabularx}
%\usepackage{makecell}


%%%%%%%%%%%%%%%%%%%%%%%%%%%%%%%%%%%%%%%%%%%%%%%%%%%%%%%%%%%%%%%%%
%%%%%%%%%%%%%%%%%%%%%%%%%%%%%%%%%%%%%%%%%%%%%%%%%%%%%%%%%%%%%%%%%



%%%%%%%%%%%%%%%%%%%%%%%%%%%%%%%%%%%%%%%%%%%%%%%%%%%%%%%%%%%%%%%%%
%%%%%%      Definice informací o dokumentu             %%%%%%%%%%
%%%%%%%%%%%%%%%%%%%%%%%%%%%%%%%%%%%%%%%%%%%%%%%%%%%%%%%%%%%%%%%%%

\input{nastaveni}  % do tohoto souboru doplňte údaje o sobě, druhu práce, názvu...

%%%%%%%%%%%%%%%%%%%%%%%%%%%%%%%%%%%%%%%%%%%%%%%%%%%%%%%%%%%%%%%%%%%%%%%%

%%%%%%%%%%%%%%%%%%%%%%%%%%%%%%%%%%%%%%%%%%%%%%%%%%%%%%%%%%%%%%%%%%%%%%%%
%%%%%%     Nastavení polí ve Vlastnostech dokumentu PDF      %%%%%%%%%%%
%%%%%%%%%%%%%%%%%%%%%%%%%%%%%%%%%%%%%%%%%%%%%%%%%%%%%%%%%%%%%%%%%%%%%%%%
%% Při načteném balíčku 'hyperref' lze použít příkaz '\pdfsettings':
\pdfsettings
%  Nastavení polí je možné provést také ručně příkazem:
%\hypersetup{
%  pdftitle={Název studentské práce},    	% Pole 'Document Title'
%  pdfauthor={Autor studenstké práce},   	% Pole 'Author'
%  pdfsubject={Typ práce}, 				    % Pole 'Subject'
%  pdfkeywords={Klíčová slova}           	% Pole 'Keywords'
%}
%%%%%%%%%%%%%%%%%%%%%%%%%%%%%%%%%%%%%%%%%%%%%%%%%%%%%%%%%%%%%%%%%%%%%%%



%%%%%%%%%%%%%%%%%%%%%%%%%%%%%%%%%%%%%%%%%%%%%%%%%%%%%%%%%%%%%%%%%%%%%%%
%%%%%%%%%%%       Začátek dokumentu               %%%%%%%%%%%%%%%%%%%%%
%%%%%%%%%%%%%%%%%%%%%%%%%%%%%%%%%%%%%%%%%%%%%%%%%%%%%%%%%%%%%%%%%%%%%%%
\begin{document}
\pagestyle{empty} %vypnutí číslování stránek

%%% Vložení desek -- od září 2021 na žádost fakulty nepoužíváno
%\includepdf[pages=1]%  buďto generovaných informačním systémem
  %{pdf/student-desky}% název souboru nesmí obsahovat mezery!
%%% NEBO vytvoření desek z balíčku
%%\makecover
%%%
%\oddpage % při dvojstranném tisku přidá prázdnou stránku
%% kazdopádně ale:
%\setcounter{page}{1} %resetovaní čítače stránek -- desky do číslování nezahrnujeme

%% Vložení titulního listu
\includepdf[pages=1]%    buďto generovaného informačním systémem
  {pdf/student-titulka}% název souboru nesmí obsahovat mezery!
%% NEBO vytvoření titulní stránky z balíčku
%\maketitle 
%%
\oddpage  % při dvojstranném tisku se přidá prázdná stránka

%% Vložení zadání
\includepdf[pages=1]%   buďto generovaného informačním systémem
  {pdf/student-zadani}% název souboru nesmí obsahovat mezery!
%% NEBO lze vytvořit prázdný list příkazem ze šablony
% \patternpage{}%
% 	{\sffamily\Huge\centering ZDE VLOŽIT LIST ZADÁNÍ}%
% 	{\sffamily\centering Z~důvodu správného číslování stránek}
%%
\oddpage% při dvojstranném tisku se přidá prázdná stránka

%% Vysázení stránky s abstraktem
\makeabstract

% Vysázení stránky s rozšířeným abstraktem
% (pokud píšete práci v češtině či slovenštině, vložení rozšířeného abstraktu zrušte;
%  pro semestrální projekt také není potřeba rozšířený abstrakt uvádět)
\input{text/rozsireny_abstrakt}

%%% Vysázení citace práce
\makecitation

%%% Vysázení prohlášení o samostatnosti
\makedeclaration

%%% Vysázení poděkování
\makeacknowledgement

%%% Vysázení obsahu
\tableofcontents

%%% Vysázení seznamu obrázků
% (vynechejte, pokud máte dva nebo méně obrázků)
\listoffigures

%%% Vysázení seznamu tabulek
% (vynechejte, pokud máte dvě nebo méně tabulek)
\listoftables

%%% Vysázení seznamu algoritmů
% (vynechejte, pokud máte dva nebo méně algoritmů)
\listofalgorithms

%%% Vysázení seznamu výpisů kódu
% (vynechejte, pokud máte dva nebo méně výpisů)
\lstlistoflistings

\cleardoublepage\pagestyle{plain}   % zapnutí číslování stránek


%%%%%%%%%%%%%%%%%%%%%%%%%%%%%%%%%%%%%%%%%%%%%%%%%%%%%%%%%%%%%%%%%
%%%%%%%%%%%%%%%%%%%       Kapitoly        %%%%%%%%%%%%%%%%%%%%%%%
%%%%%%%%%%%%%%%%%%%%%%%%%%%%%%%%%%%%%%%%%%%%%%%%%%%%%%%%%%%%%%%%%

%Pro vkládání kapitol i příloh používejte raději \include než \input
%%% Vložení souboru 'text/uvod.tex' s úvodem
\include{text/uvod}

%%% Vložení souboru 'text/cile.tex' s úvodem
\include{text/cile}

%%% Vložení souboru 'text/reseni' s popisem řešení práce
% (rozdělte na více souborů či kapitol, pokud je vhodné)
\include{text/reseni}

%%% Vložení souboru 'text/vysledky' s popisem vysledků práce
% (rozdělte na více souborů či kapitol, pokud je vhodné)
\include{text/vysledky}

%%% Vložení souboru 'text/zaver' se závěrem
\include{text/zaver}


%%%%%%%%%%%%%%%%%%%%%%%%%%%%%%%%%%%%%%%%%%%%%%%%%%%%%%%%%%%%%%%%%
%%%%%%%%%%%%%%%%      Tisk literatury        %%%%%%%%%%%%%%%%%%%%
%%%%%%%%%%%%%%%%%%%%%%%%%%%%%%%%%%%%%%%%%%%%%%%%%%%%%%%%%%%%%%%%%

%%% Při využití vlastní definice pomocí bibitem
%%% Vložení souboru 'text/literatura' se seznamem zdrojů
\include{text/literatura}


%%% Při využití literatury pomocí BIBLATEX
% === běžná literatura ===
%\printbibliography[
%  heading=bibintoc,
%  notkeyword={P},
%  title={Literatura}
%]

%%% Pro disertační práce, oddělená literatura autora.
%%% Každý takový zdroj musí obsahovat položku: keyword={P}
% === publikace autora ===
%\newrefcontext[labelprefix=P]
%\printbibliography[
%  heading=bibintoc,
%  keyword={P},
%  title={Literatura autora}
%]


%%%%%%%%%%%%%%%%%%%%%%%%%%%%%%%%%%%%%%%%%%%%%%%%%%%%%%%%%%%%%%%%%
%%%%%%%%%%%%%%%%    Tisk zkratek, příloh     %%%%%%%%%%%%%%%%%%%%
%%%%%%%%%%%%%%%%%%%%%%%%%%%%%%%%%%%%%%%%%%%%%%%%%%%%%%%%%%%%%%%%%

%%% Vložení souboru 'text/zkratky' se seznam použitých symbolů, veličin a zkratek
\include{text/zkratky}

%%% Začátek příloh
\appendix

%%% Vysázení seznamu příloh
% (vynechejte, pokud máte dvě nebo méně příloh)
\listofappendices

%%% Vložení souboru 'text/prilohy' s přílohami
% Obvykle je přítomen alespoň popis co najdeme na přiloženém médiu
\include{text/prilohy}

\end{document}
